% Part: first-order-logic
% Chapter: proof-systems
% Section: tableaux

\documentclass[../../../include/open-logic-section]{subfiles}

\begin{document}

\iftag{FOL}
      {\olfileid{fol}{prf}{tab}}
      {\olfileid{pl}{prf}{tab}}

\olsection{\usetoken{P}{tableau}}

د !!{derivation} ډېر نظامونه د !!{sentence}s له ترتيبونو سره کار
کوي، خو !!{tableau}s له !!{signed formula}s سره کار کوي۔
!!^a{signed formula} يوه جوړه ده چې د رښتياوالي له يوې نښې
($\True$ يا $\False$) او !!a{sentence} څخه جوړه ده:
\[
\sFmla{\True}{!A} \text{ يا } \sFmla{\False}{!A}.
\]
!!^a{tableau} له هغو !!{signed formula}s څخه جوړ وي چې د لاندې خوا
ته څانګېدونکې ونه کښې ترتيب شوي وي۔ دا له يو شمېر \emph{فرضونو}
پيلېږي او په هغو !!{signed formula}s دوام کوي چې د استنتاج د کومې
قاعدې په پلي کولو تر ځان له پاسه له يوه !!{signed formula}s څخه
راوتلي وي۔ هره قاعده اجازه راکوي چې د يوې څانګې په پاے کښې يو يا
څو !!{signed formula}s ورزيات کړو، يا دوه !!{signed formula}s يو د
بل تر څنګ ورزيات کړو؛ په وروستي حالت کښې څانګه په دوو وېشل کېږي او
دوه ورزيات شوي !!{signed
  formula}s د دواړو څانګو سرونه جوړوي۔

پر يوه مرکب !!{signed formula} د قاعدې له پلي کولو څخه داسې
!!{signed formula}s ورزياتېږي چې سمدستي فرعي-!!{formula}s وي۔ د هرې
قاعدې يوه جوړه راځي، د دواړو نښو هرې يوې دپاره يوه قاعده۔ مثلاً د
$\TRule{\True}{\land}$ قاعده پر $\sFmla{\True}{!A \land !B}$ پلې
کېږي او اجازه راکوي چې د هرې هغې څانګې په پاے کښې دواړه
!!{signed formula}s $\sFmla{\True}{!A}$ او~$\sFmla{\True}{!B}$
ورزيات کړو چې $\sFmla{\True}{!A \land !B}$ لري؛ او د
$\TRule{\False}{\land}$ قاعده اجازه راکوي چې يوه څانګه د
$\sFmla{\False}{!A}$ او $\sFmla{\False}{!B}$ يو د بل تر څنګ په
ورزياتولو ووېشو۔
\textbf{د ژباړې د نسخې سمون (OLPF-002):} په سرچينه کښې د دروغې
عطف قاعدې دويم ارګومېنټ ټول مرکب فارمول دے، خو د رښتوني عطف موازي
قاعده او د لاندې بېلګې ټولې قاعدې هلته يوازې نښلوونکے اخلي۔ دلته
هماغه د عطف نښلوونکے کارول شوے، څو د قاعدې نښه په يو شان بڼه جوړه
شي۔
!!^a{tableau} هله تړلے دے چې د هغه هره څانګه د !!{signed formula}s
سره برابره جوړه $\sFmla{\True}{!A}$ او $\sFmla{\False}{!A}$ ولري۔

پر !!{tableau}s ولاړه $\Proves$ اړيکه داسې تعريفېږي:
$\Gamma \Proves !A$ هله او يوازې هله چې داسې کوم متناهي سټ
~$\Gamma_0 = \{!B_1, \dots, !B_n\} \subseteq \Gamma$ موجود وي چې د
لاندې فرضونو دپاره يو تړلے !!{tableau} موجود وي:
\[
\{\sFmla{\False}{!A}, \sFmla{\True}{!B_1}, \dots, \sFmla{\True}{!B_n}\} 
\]
مثلاً دا تړلے !!{tableau} ښيي چې $\Proves (!A \land !B) \lif !A$:
\begin{oltableau}
  [\sFmla{\False}{(\formula{A} \land \formula{B}) \lif \formula{A}}, just = \TAss
    [\sFmla{\True}{\formula{A} \land \formula{B}}, just = {\TRule{\False}{\lif}[1]}
      [\sFmla{\False}{\formula{A}}, just = {\TRule{\False}{\lif}[1]}
        [\sFmla{\True}{\formula{A}}, just = {\TRule{\True}{\lif}[2]}
          [\sFmla{\True}{\formula{B}}, just = {\TRule{\True}{\lif}[2]}, close]
        ]
      ]
    ]
  ]
\end{oltableau}

يو سټ $\Gamma$ د !!{tableau} په حساب کښې هله ناسازګار دے چې د
لاندې فرضونو دپاره يو تړلے !!{tableau} موجود وي:
\[
\{\sFmla{\True}{!B_1}, \dots, \sFmla{\True}{!B_n}\} 
\]
د کوم $!B_i \in \Gamma$ دپاره۔

!!^{tableau}s په ۱۹۵۰مو کلونو کښې اېورټ بېت او ياکو هېنتيکا هر يوه
په خپلواکه توګه جوړ کړل، او رېمونډ سموليان ساده او عام کړل۔ دا کارول
ډېر اسانه دي، ځکه چې د !!a{tableau} جوړول ډېره منظمه طريقه ده۔ د
!!{tableau}s د منظم طبيعت له امله په کمپيوټر کښې هم په اسانۍ پلي
کېږي۔ خو !!a{tableau} زياتره په سختۍ لوستل کېږي او له ثبوتونو سره ئې
تړاو کله کله په اسانۍ نۀ ښکاري۔ دا تګلاره هم ډېره عامه ده، او ډېر
بېل منطقونه د !!{tableau} نظامونه لري۔ !!^{tableau}s له موږ سره د
هغو !!{structure}s په موندلو کښې هم مرسته کوي چې ورکړل شوي
!!{sentence}s، يا د هغو سټونه، صادقوي: که سټ د صدق وړ وي، تړلے
!!{tableau} به و نۀ لري؛ يعنې هر !!{tableau} به يوه پرانيستې څانګه
ولري۔ په دې شرط چې پر هغې څانګه هره ممکنه قاعده پلې شوې وي،
صادقوونکے !!{structure} له يوې پرانيستې څانګې څخه ``لوستل کېدے''
شي۔ له سېکوېنټ حساب سره هم ډېر نږدې تړاو شته: په اصل کښې يو تړلے
!!{tableau} د سېکوېنټ په حساب کښې يو لنډ شوے !!{derivation} دے چې
سرچپه ليکل شوے وي۔

\end{document}
